% LSPV-2 REPORT::INFOFLOWCHAIN1::2026-10-12T17:27:40Z
% !TEX program = xelatex
\documentclass[11pt]{article}
\usepackage[margin=0.88in]{geometry}
\usepackage{fontspec,amsmath,amssymb,amsthm,booktabs,longtable,array,enumitem,xcolor,hyperref,xurl,microtype}
\setmainfont{Latin Modern Roman}
\hypersetup{colorlinks=true,linkcolor=blue!50!black,urlcolor=blue!60!black,citecolor=blue!50!black}
\newtheorem{theorem}{Theorem}
\newtheorem{corollary}[theorem]{Corollary}
\newcommand{\EXACT}{\textbf{PROVED/EXACT}}
\newcommand{\KNOWN}{\textbf{KNOWN THEOREM}}
\newcommand{\COMP}{\textbf{COMPUTATIONAL}}
\newcommand{\CAND}{\textbf{CANDIDATE}}
\newcommand{\OPEN}{\textbf{OPEN}}
\newcommand{\RETIRED}{\textbf{FALSIFIED/RETIRED}}
\newcommand{\claim}[1]{\texttt{#1}}

\title{LSPV-2 universal-core test: temporal information gain is source transport, not semantic work}
\author{UNIVERSAL\_CORE\_TESTS}
\date{2026-10-13 02:27 KST}

\begin{document}
\maketitle

\section*{Executive status}
\textbf{P versus NP remains \OPEN.}  This run tests the most immediate repair
to INFONOGO1: replace aggregate mutual information by transition-by-transition
conditional information.  The exact chain rule shows that this repair does
not distinguish compatibility computation from invertible source transport.
The theorem \claim{TEMPORAL-INFORMATION-TRANSPORT-TRILEMMA1} is \EXACT,
pending independent review.  The implication
\claim{PER-STEP-INFORMATION-GAIN-AS-SEMANTIC-COST1} is \RETIRED.  A
machine-defined, representation-invariant, same-execution computational
provenance increment remains \OPEN.

\section{Predecessor write-recovery preflight}
The immediately preceding LSPV-1 run is present at
\path{research/agents/lspv-1/2026-10-13_0209_LSPV-1_INFONOGO1.tex}.  GitHub
readback found the required header
\texttt{LSPV-1 REPORT::INFONOGO1} and final verification sentinel.  The report
is substantive, so no recovery request or duplicate file was needed.

\section{Barrier and primary-literature preflight}
The two blocking obstacles are exact.
\begin{enumerate}[leftmargin=*]
\item A numerical information-flow functional must distinguish mere copying,
invertible recoding, delay and replay from performing the chosen semantic
operation $H$.
\item Even if such a distinction existed, a theorem must convert it into time
or space of the \emph{same ordinary uniform deciding execution}, with encoded
length and every observation/reconstruction cost charged, and then land in an
unconditional lower bound for that same model.
\end{enumerate}
Shannon's entropy formalism supplies the chain rule used below
\cite{shannon1948}.  Massey's directed information sums conditional mutual
information along a causal order \cite{massey1990}.  Schreiber's transfer
entropy was introduced precisely because delayed mutual information can
confound actual exchange with common history or common input
\cite{schreiber2000}.  These are primary information-theory sources.  None
states that positive information flow is computational work on a deterministic
read-only input.  Baker--Gill--Solovay, Razborov--Rudich,
Aaronson--Wigderson, and Fortnow--Lipton--van Melkebeek--Viglas delimit the
complexity-theoretic transfer problem \cite{bgs1975,rr1997,aw2009,flmv2005}.

\paragraph{Transfer classification.}
The Shannon chain rule is a \textbf{directly transferable theorem}.  Transfer
entropy's conditioning lesson is a \textbf{useful structural analogy requiring
a new transfer theorem}: conditioning on all common source input makes the
deterministic quantity zero, not costly.  Interpreting a positive temporal
information total as ordinary-machine time or space is a \textbf{misleading
analogy / failed transfer}.  VLSI routing and communication-cut lower bounds
remain restricted-model results unless an ordinary SAT computation is reduced
to the same physical or communication interface with all costs preserved.

\section{Exact quantifiers, encoding, and candidate}
Fix an encoded input length $N$ and a finite length layer $\mathcal X_N
\subseteq\{0,1\}^N$ with any distribution (uniform in the audit).  Let
\[
 A=a(X)\in\{0,1\},\qquad H=h(X),\qquad
 S_{1:m}=(S_1,\ldots,S_m)=s(X)
\]
be deterministic functions.  Here $A$ is the decision bit, $H$ is any chosen
source-defined compatibility object, and $S_t$ is an observed event or state
slice from one run.  No claim is made that $S_t$ is a canonical machine state.
Define the directed-information-style static-source increments
\[
 g_t:=I(H;S_t\mid A,S_{<t}).
\]
This is the static-source specialization relevant to the proposed repair; it
is not a claim that every standard definition of Massey directed information
is identical to this expression.

The candidate implication under attack was
\[
 \sum_{t=1}^{m}g_t>0\quad\text{or large}
 \quad\Longrightarrow\quad
 \text{the deciding execution paid for $H$-specific semantic work}.
 \tag{C}
\]
Any P-versus-NP-strength form must hold for every uniform deterministic SAT
decider, every allowed equivalent wrapper, and a worst-case input depending on
the decider, with all quantities measured in the encoded length $N$.

\section{The temporal information transport theorem}
\begin{theorem}[\claim{TEMPORAL-INFORMATION-TRANSPORT-TRILEMMA1}]
For the finite random variables above,
\begin{equation}
 \sum_{t=1}^{m} I(H;S_t\mid A,S_{<t})
   = I(H;S_{1:m}\mid A)
   \le H(H\mid A)\le H(X\mid A)\le N.       \tag{1}
\end{equation}
If $X$ is a deterministic function of $(A,S_{1:m})$, then equality holds in
the first upper bound:
\begin{equation}
 \sum_{t=1}^{m} I(H;S_t\mid A,S_{<t})=H(H\mid A). \tag{2}
\end{equation}
In contrast, if the common source $X$ is included in every condition, then
\begin{equation}
 \sum_{t=1}^{m}I(H;S_t\mid X,A,S_{<t})=0.       \tag{3}
\end{equation}
Thus the unconditioned temporal score is maximized by lossless transport of
the source, while conditioning away that transport annihilates every
deterministic source-defined target.
\end{theorem}

\begin{proof}
The conditional mutual-information chain rule gives the equality in (1).
The first inequality is $I(U;V\mid W)\le H(U\mid W)$.  Since $H=h(X)$,
conditional data processing gives $H(H\mid A)\le H(X\mid A)$, and an
$N$-bit source has entropy at most $N$.  If $(A,S_{1:m})$ determines $X$, it
determines $H$, so $H(H\mid A,S_{1:m})=0$ and (2) follows.  Finally, $H$ is
already determined by $X$, hence every term in (3) is zero.
\end{proof}

\begin{corollary}[serial COPY and prefix-parity transport]
Let $S_t=X_t$ for $t\le N$.  Alternatively let
$S_t=X_1\oplus\cdots\oplus X_t$.  In both cases $S_{1:N}$ determines $X$,
so the temporal score equals $H(H\mid A)$ for every deterministic target
$H=h(X)$.  The first stream is source COPY.  The second is an invertible
online linear recoding, because $X_t=S_{t-1}\oplus S_t$ with $S_0=0$.
Neither stream need evaluate $h$.
\end{corollary}

\section{Explicit forced-assignment 3CNF control}
For $b=(b_1,b_2,b_3)\in\{0,1\}^3$, construct an equal-length 3CNF $F_b$.
For each target variable $z_i$, put $\ell_i=z_i$ when $b_i=1$ and
$\ell_i=\neg z_i$ when $b_i=0$, introduce fresh $u_i,v_i$, and include
\[
(\ell_i\vee u_i\vee v_i),\ (\ell_i\vee u_i\vee\neg v_i),\
(\ell_i\vee\neg u_i\vee v_i),\ (\ell_i\vee\neg u_i\vee\neg v_i).
\]
The four clauses force $\ell_i=1$ and otherwise leave $u_i,v_i$ free.  Hence
all eight formulas are YES instances of equal encoding length and
$H(F_b)=b$ is the forced target assignment while $A(F_b)=1$.  Serial COPY or
prefix-parity transport of the three label bits has increments $(1,1,1)$ and
total $3$ bits; DEAD and answer-only views have total $0$.  The construction
is a width-three forced-assignment control, not a claim that it meets every
additional regular signed clean-core promise.

\section{All-cost and SAME\_REALIZATION audit}
A read-only-input wrapper can emit $(t,X_t)$ serially before simulating its
base decider.  It adds $N$ source reads, $N$ writes, $O(N)$ time and
$\lceil\log_2(N+1)\rceil+O(1)$ auxiliary bits.  Prefix parity adds $N$ XOR
updates with the same asymptotic bounds.  Reconstructing the source from the
recorded trace costs $N$ reads/writes (and $N-1$ XORs for prefix parity).
Evaluating $h$ after reconstruction is a separate, fully charged operation.

The transport occurs in the same wrapped execution, so this is not a
different-realization objection.  What fails is the semantic inference:
maximum temporal information can be present although the deciding run never
evaluates $h$.  An external recorder may reconstruct $X$ and evaluate $h$ or
call a solver again, but that is external extraction, not proof that the first
run paid for $H$-specific work.  Equation (1) additionally caps every such
source-derived Shannon total by $N$ bits, so the functional alone cannot yield
a superlinear, let alone superpolynomial, time lower bound.

\section{Bounded reproducible hostile-control audit}
The script
\path{experiments/review/per_step_information_copy_audit.py} enumerated all
binary sources for $1\le N\le8$, six targets, four decision quotients, and
eight streams: serial COPY, prefix parity, reverse COPY, DEAD, answer-only,
DELAY, INTERP and RECOMP.  It checked (1), (3), lossless-stream equality and
collapsed-stream zero, then checked the eight-formula forced-assignment
control.  The run passed 4,614 assertions in 6.566457 seconds.  Results are in
\path{results/raw/review/2026-10-13_per_step_information_copy.json}.

The ledger explicitly includes 2SAT, Horn, affine/Tseitin, bounded-treewidth,
independent-product and signed-clean-core controls; padding, redundant
clauses, literal renaming and clause reordering; five encodings; and
DIRECT/DEAD/COPY/DELAY/INTERP/RECOMP/PREFIX-PARITY wrappers.  These labels
check that the exact information identity is representation-agnostic.  They
do not establish asymptotic SAT cost or claim every label instance realizes
the forced-assignment gadget.  Finite results are \COMP{} only; the theorem,
not enumeration, supplies the all-$N$ statement.  No invocation exceeded the
30-second runtime rule and no chunking was required.

\section{Counterexample and smallest admissible repair}
Candidate (C) is \RETIRED{} because two cheap transports attain its maximum.
The smallest syntactic repair is to quotient out lossless source transport or
condition on the common input.  The theorem shows the dilemma: conditioning
on all of $X$ makes every deterministic $H=h(X)$ increment zero, whereas a
weaker quotient must define which recodings, delays, recomputations and
wrappers count as semantically free.  Such a quotient is a representation or
provenance choice, not yet an invariant of arbitrary computation.

The surviving target is therefore:
\begin{quote}
\claim{COMPUTATIONAL-PROVENANCE-INCREMENT1} (\OPEN): define a uniform,
machine-model-fixed increment that is invariant under lossless COPY, invertible
online recoding, delay, interpolation, replay, recomputation and alternate
encodings, is positive only when the same deciding execution performs an
$H$-typed operation, lower-bounds a charged ordinary-machine resource, and
lands in a model-matched unconditional forbidden region.
\end{quote}
No such definition or theorem is supplied here.

\section{Relativization, algebrization, Natural Proofs, and P=NP audit}
The negative theorem relativizes: an oracle machine can be prefixed by the
same serial or parity transport.  Its proof is an entropy identity and also
survives algebraic extensions; it is not the required nonalgebrizing hinge.
It defines no large constructive property of Boolean functions, so it does
not evade or invoke the Natural Proofs barrier.  Under $P=NP$, substitute any
polynomial SAT decider into the wrapper.  The total remains polynomial and no
known unconditional ordinary-machine theorem forbids it.  Promoting the open
provenance increment to a superpolynomial cost would simply restate the
missing lower bound.

\section{Post-quantum technique ledger}
\begin{longtable}{p{0.28\linewidth}p{0.64\linewidth}}
\toprule
Technique & Classification and audit\\ \midrule
Lossiness/collision entropy & \textbf{EXACT OBJECT/THEOREM TRANSFER}.  COPY and prefix parity are injective length-layer maps, so their collision count is $|\mathcal X_N|$ and their images retain the full source entropy.  This is algorithmically exploitable transport, not hardness.\\
Schur squares; quadratic relation spaces; Hilbert functions and syzygies & \textbf{FAILED TRANSFER} for the current barrier.  They are exact invariants only after choosing a code/evaluation representation and do not distinguish semantic evaluation from streaming that representation.\\
Worst-case/average-case black-box extraction & \textbf{ANALOGY/CANDIDATE needing a new bridge}.  External calls on restrictions have the right solver-independent flavor, but exponent preservation, average-case validity and SAME\_REALIZATION are open.\\
Collision extraction / special soundness & \textbf{FAILED TRANSFER}.  Two source encodings are different statements, not conflicting openings of one commitment, so no forbidden witness follows.\\
Hash/random oracle methods & \textbf{FAILED TRANSFER} for an unconditional separation unless every hardness and oracle assumption is removed.\\
\bottomrule
\end{longtable}

\section{CORPUS REVIEWED}
The repository tree was surveyed before substantive work.  The review read
root \path{main.tex}, \path{AGENTS.md}, \path{CODEX.md}, the complete
\path{.codex/skills/pnp-research/SKILL.md}, actual \path{main4.tex} and
\path{main5.tex}, every task specification under \path{tasks/}, and the newest
verified MAIN \path{research/main/2026-10-13_0148_MAIN_TRACECONJ1.tex}.
Material peer inputs were INFONOGO1, SERIALTRACE1, CAUSALCANCEL1 and
STATE\_READBACK1.  Relevant older work surfaced by repository search included
the Lipschitz-squeeze information audit, PRODSTREAM1, DERIVEDLEAK1,
JOINTENDPOINT1 and the HEC/VSC-HEC/G-PPC-HEC/CRS-PPC/CCSPC/RCCN/Search-G
causal branch.  The governance reconciliation note, hostile-control tasks and
the active experiment/result artifacts were inspected.  Repository search for
``directed information'', ``transfer entropy'', ``conditional mutual
information'', and ``per-step information'' found no prior theorem stating
the exact transport trilemma proved here.

\section{DUPLICATION / CONTRADICTION AUDIT}
INFONOGO1 already proved that aggregate $I(H;T\mid A)$ is either decision-only,
zero after conditioning on $X$, or copyable.  The older Lipschitz audit noted
that Shannon information originating in an $N$-bit source has total mass at
most $N$.  SERIALTRACE1 proved source injection through a serial full trace,
and CAUSALCANCEL1 retired bare path liveness.  The new contribution is the
exact temporal localization: the sum of all stepwise gains is precisely the
same aggregate quantity, and prefix parity shows that a nonliteral online
recoding also saturates it.  Thus this report strengthens and closes the
suggested directed-information repair; it does not rename a prior result or
contradict any verified theorem.

\section{WHAT CHANGED RELATIVE TO ROOT MAIN}
The LSPV-2 ledger now records
\claim{TEMPORAL-INFORMATION-TRANSPORT-TRILEMMA1} as \EXACT,
\claim{PER-STEP-INFORMATION-GAIN-AS-SEMANTIC-COST1} as \RETIRED, and
\claim{COMPUTATIONAL-PROVENANCE-INCREMENT1} as \OPEN.  No executive theorem,
SAT exponent, endpoint, or claim in another specialist's ownership region is
strengthened.

\section{ADOPTED / REJECTED / RETIRED PRIOR COMPONENTS}
\begin{itemize}[leftmargin=*]
\item \textbf{Adopted:} INFONOGO1's deterministic-target trilemma,
SERIALTRACE1's serial source injection, CAUSALCANCEL1's representation warning,
and the restricted HEC/RCCN results in their stated models.
\item \textbf{Rejected as a universal bridge:} raw temporal conditional
mutual information, including positive per-step gain and maximum total gain.
\item \textbf{Retained \OPEN:} a wrapper-invariant computational-provenance
increment with an exact solver-to-resource and resource-to-forbidden-endpoint
transfer.
\end{itemize}

\section{STALE OR OVERSTATED CLAIMS FOUND}
The pre-run root correctly called directed information a candidate, not a
theorem.  After this audit, the static-source per-step conditional-information
version is stale as a live candidate and must be retired in that scope.
Standard directed information or transfer entropy should not be described as
computational lower bounds; their information-flow meanings do not include the
required ordinary-machine landing.  The HEC/RCCN causal theorems remain valid
only for their explicit online/network services and must not be silently
promoted to arbitrary static SAT deciders.

\section{Faithful comparison with actual main4 and main5}
Actual \path{main4.tex} is the ARC achievable-region manuscript and explicitly
leaves its SAT-specific finite-hitting/charged-successor premise open.  Actual
\path{main5.tex} is the WRAPPER manuscript; it fixes a supplied polynomial
decider exponent and leaves a favorable charged interface open.  Neither
document asserts this temporal-information theorem, nor does either claim
that semantic-state count implies runtime.  This report changes neither
manuscript and does not substitute the present negative theorem for their open
constructive interfaces.

\section{Handoff and exact missing arrow}
Future synthesis should reject any claimed semantic charge based solely on
\(\sum_t I(H;S_t\mid A,S_{<t})\), on its positivity, or on saturation of
\(H(H\mid A)\).  A new proposal must provide an operational primitive that
cannot be implemented by lossless transport or invertible streaming recoding,
prove invariance under equivalent wrappers and encodings, charge observation
and validation, and lower-bound time/space of the same ordinary execution.
The final model-matched forbidden endpoint is separately \OPEN.  Therefore
P versus NP remains \OPEN.

\begin{thebibliography}{9}
\bibitem{shannon1948} C. E. Shannon, ``A Mathematical Theory of Communication,'' \emph{Bell System Technical Journal} 27 (1948), 379--423, 623--656. \url{https://people.math.harvard.edu/~ctm/home/text/others/shannon/entropy/entropy.pdf}
\bibitem{massey1990} J. L. Massey, ``Causality, Feedback and Directed Information,'' \emph{Proc. ISITA}, 1990, 303--305.
\bibitem{schreiber2000} T. Schreiber, ``Measuring Information Transfer,'' \emph{Physical Review Letters} 85 (2000), 461--464. \url{https://arxiv.org/abs/nlin/0001042}
\bibitem{bgs1975} T. Baker, J. Gill, and R. Solovay, ``Relativizations of the P=?NP Question,'' \emph{SIAM J. Comput.} 4 (1975), 431--442. \url{https://doi.org/10.1137/0204037}
\bibitem{rr1997} A. Razborov and S. Rudich, ``Natural Proofs,'' \emph{J. Comput. Syst. Sci.} 55 (1997), 24--35. \url{https://doi.org/10.1006/jcss.1997.1494}
\bibitem{aw2009} S. Aaronson and A. Wigderson, ``Algebrization: A New Barrier in Complexity Theory,'' \emph{ACM Trans. Comput. Theory} 1 (2009). \url{https://doi.org/10.1145/1490270.1490272}
\bibitem{flmv2005} L. Fortnow, R. Lipton, D. van Melkebeek, and A. Viglas, ``Time-Space Lower Bounds for Satisfiability,'' \emph{J. ACM} 52 (2005), 835--865. \url{https://doi.org/10.1145/1101821.1101822}
\end{thebibliography}

% END_OF_VERIFIED_REPORT::LSPV-2::INFOFLOWCHAIN1
\end{document}
